\section{引言}
\label{sec:intro}

变分量子本征求解器结合参数化电路与经典优化，以近似低能态~\cite{peruzzo2014variational}。
超导处理器实验已为分子和量子磁体建立硬件高效实现~\cite{kandala2017hardware}，后续综述系统讨论了电路设计、测量成本和优化的作用~\cite{cerezo2021variational,tilly2022variational}。
对于拓扑物态，保护对称性与非局域关联共同指导电路设计和观测量选择。

一维对称有能隙自旋相可通过对称性在纠缠空间中的投影作用区分~\cite{chen2011classification,schuch2011classifying}。
这种 SPT 区分与内禀拓扑序不同，后者可用拓扑纠缠熵诊断~\cite{kitaev2006topological,levin2006detecting}。
超导处理器已在 toric-code 态中测量该量~\cite{satzinger2021realizing}。
本文研究具有平庸、Haldane 型和反铁磁（AFM）竞争区域的相互作用二聚化自旋链。

里德伯原子实验实现了玻色 SSH 的 SPT 相~\cite{deleseleuc2019observation}。
数字实验已在具有次近邻自旋交换的模型中制备 SPT 相~\cite{tan2023realizing}，并利用云端超导设备模拟更长的自旋链~\cite{yu2023simulating}。
在变分制备方面，Sun 等将结构化初态与变分层结合，研究关联拓扑相~\cite{sun2023efficient}。
这些工作启发我们考察：当平庸二聚化、拓扑二聚化与 AFM 关联竞争时，依据模型选择的初态和电路能达到何种效果。

测量选择决定了能够从制备态推断哪些性质。
部分反射不变量提供基于对称性的多体诊断~\cite{elben2020manybody}，普通 string 序的解释则依赖串内对称操作与端点算符~\cite{pollmann2012detection,haegeman2012order}。
本文以部分反射不变量构造数值相图，在实验中以单一计算基测量交错关联和晶胞 string 关联。
三个参考态确定这两类关联量的尺度及其仿射组合 $Q$。

本文以相匹配初态和反射对称变分电路为三种区域提供共同框架。
在理想经典模拟中优化电路参数后，于百花处理器执行优化后的态制备电路并测量关联诊断量。
数值相图取 $L=8$；另一组固定 $\delta=-2$ 的对称性扇区能隙计算比较五种链长。
读出缓解采用综述~\cite{endo2021hybrid,cai2023quantum} 所述的分配矩阵框架。
实验分析比较 $30$ 个扫描点上的 $Q$，并给出拓扑区的正晶胞 string 关联。

下文依次介绍模型的对称性分类、有限尺寸相图、测量诊断量、电路构造、读出缓解及实验与数值模拟对比。
附录给出体理论背景和有限尺寸能隙分析，再介绍比特筛选结果。
